\subsection{A functorial reduction of both integral uniqueness questions}

For a collection of sets $E=(E_w)_{w\geq1}$, let $P(E)$ be the
one-truncated PD DGA defined by (SK.1), with one disk
$x_{w,e}\mapsto y_{w,e}$ for each $e\in E_w$ and the same
endpoint values $0,T^{[w]}$. Empty and infinite sets are allowed.
Let $\mathcal C\subset\mathcal A_{\mathbb Z}$ be the full
subcategory of objects whose first skeleton admits a strict
isomorphism with some $P(E)$, where $E_w$ is nonempty
for every sufficiently large $w$.
The isomorphism is a property of the object and is not additional
morphism data; the functor below supplies its own explicit
identification. Write $j:\mathcal C\hookrightarrow\mathcal A_{\mathbb Z}$
for the full-subcategory inclusion.
All sets and categories are taken in a fixed sufficiently large
universe, preserved by the indicated free-algebra constructions.

\begin{theorem}[A-006: Proved; a strict functorial skeleton replacement]
\label{A006}
There is an explicit functor
$R:\mathcal A_{\mathbb Z}\to\mathcal C$ and an actual natural
transformation of strict PD-DGA maps
\[
                         \tau:jR\Longrightarrow
                                   \mathrm{Id}_{\mathcal A_{\mathbb Z}}.
 \tag{SK.15}
\]
Its restriction to the full subcategory is a natural transformation
$Rj\Rightarrow\mathrm{Id}_{\mathcal C}$.
Consequently inclusion and replacement induce mutually inverse
equivalences on the groupoids obtained by inverting every morphism,
and their maps of nerves are homotopy inverses.

Common admissible sources for all pairs in $\mathcal A_{\mathbb Z}$
are equivalent to common sources internal to $\mathcal C$ for all
pairs in $\mathcal C$.
Terminality of the all-arrow localization of
$\mathcal A_{\mathbb Z}$ is equivalent to terminality of the
all-arrow localization of $\mathcal C$.
All statements remain valid in the support-restricted category,
the termwise-flat integral category, and their intersection,
with the corresponding full subcategory $\mathcal C$.
\end{theorem}

\begin{proof}
For an actual admissible object $A$, define the set
\[
 E_w(A)=\{a\in A_1^{w,0}:d_1a=0,\quad d_0a=T^{[w]}\},
                    \qquad w\geq1.                              \tag{SK.16}
\]
For $w\geq\max\{1,\mu_A(0)\}$ it is nonempty by normalized
row exactness, exactly as in (SK.7).
Set $P_A=P(E(A))$. There is a canonical strict truncated map
\[
 p_A:P_A\to A_{\leq1},
             \qquad x_{w,a}\mapsto a,\quad y_{w,a}\mapsto da.
 \tag{SK.17}
\]
For a strict morphism $f:A\to B$, the element $f_1(a)$ belongs
to $E_w(B)$ because $f$ preserves faces and coefficients.
Define $P_f:P_A\to P_B$ by
\[
             x_{w,a}\mapsto x_{w,f_1(a)},\qquad
             y_{w,a}\mapsto y_{w,f_1(a)},\qquad S\mapsto S.
 \tag{SK.18}
\]
The free PD presentation gives an actual strict map, even when
different indices have the same image under $f_1$.
Identity and composition in (SK.18) are equalities on every
generator and therefore equalities of maps. Moreover
$p_BP_f=f_{\leq1}p_A$ on every generator and coefficient.
Thus $P_A$ and its evaluation $p_A$ are functorial, with no
choice of an individual edge filler or of a numerical cutoff.

Using the explicit coskeleton (SK.8), define
\[
            R(A)=A\mathop{\times}_{C(A_{\leq1})}C(P_A),
       \qquad R(f)(a,c)=(f(a),C(P_f)(c)).                         \tag{SK.19}
\]
The equality just verified makes the displayed pair a member of
the target pullback. Identities and composites follow componentwise.
Projection to the first factor is $\tau_A$; the componentwise
formula proves $f\tau_A=\tau_BR(f)$ exactly.

We verify that (SK.19) has the claimed target category. The
possibly infinite tensor product defining $(P_A)_1$ is a filtered
union of finite tensor products of the disks (SK.5).
Their integral coefficient contractions prove its vertical
coefficient comparison in every weight, without a restriction
on the cardinality of $E_w(A)$.
Eventual nonemptiness in (SK.16) proves the endpoint surjectivity
used in (SK.6), with $N_A=\max\{1,\mu_A(0)\}$.
The proof of A-004 uses precisely these two properties of $P(N)$
and the actual truncated map to $A$. Apply its displayed graph
sequences and pullback sequences to $P_A$ instead. They give
\[
 \mu_{R(A)}(d)=\max\{N_A,\mu_A(d)\},\qquad
 \lambda_{R(A)}(d)=\max\{N_A,\lambda_A(d+2),\mu_A(d+2)\}.
 \tag{SK.20}
\]
All products in each matching diagram are finite, even when
the modules being multiplied have infinitely many generators.
Their exactness and all stated uniform bounds are therefore
unchanged. The first-skeleton identifications are the maps
$c\mapsto(p_A(c),c)$ and second projection from A-004.
They commute with (SK.18)--(SK.19), so transporting through them
is also functorial. This proves $R(A)\in\mathcal C$ and proves
that every $\tau_A$ is an actual admissible morphism.
The support and flatness arguments of A-004 apply to these same
modules: the free disk terms are free abelian, their finite
matching subgroups are torsionfree, and pullbacks into a flat
target remain torsionfree.

For $C\in\mathcal C$, both $R(jC)$ and $C$ are objects of the
full subcategory $\mathcal C$. Hence the very map $\tau_C$
is a morphism in $\mathcal C$, and its established naturality
restricts there. This gives the second natural transformation.
After every morphism is inverted, these two transformations
become natural isomorphisms
\[
       (jR)^{\mathrm{loc}}\cong\mathrm{Id},\qquad
       (Rj)^{\mathrm{loc}}\cong\mathrm{Id}.
 \tag{SK.21}
\]
Thus the induced functors are inverse equivalences of the
specified localizations.

For completeness the nerve statement follows from the actual
natural transformations as well. A natural transformation
$F\Rightarrow G$ defines a functor on the product with the
two-object interval category $[1]$: use $F$ at its first object,
$G$ at its second, and the transformation on its cross arrows.
Naturality is exactly the equality ensuring composition of those
arrows. Taking nerves gives a simplicial homotopy
$N\mathcal D\times\Delta[1]\to N\mathcal E$ between $NF$ and
$NG$. Apply this construction to the two transformations above.
The nerve maps of $j$ and $R$ are therefore homotopy inverses.

We spell out the common-source equivalence because it is also
an assertion about actual arrows. Suppose common sources exist
in the full admissible category, and let $C_1,C_2\in\mathcal C$.
For a source $D\to C_1,C_2$, apply $R$ and then compose
\[
                     R(D)\longrightarrow R(C_i)
                              \xrightarrow{\tau_{C_i}} C_i.
 \tag{SK.22}
\]
This is a common source entirely inside $\mathcal C$.
Conversely, apply the internal assertion to $R(A),R(B)$;
compose its actual source maps with $\tau_A,\tau_B$ to obtain
a source for $A,B$. The optional support and flatness restrictions
transfer by the already proved properties of $R$.
Finally equivalence (SK.21) preserves and reflects equivalence
to the one-object groupoid with trivial automorphism group.
This proves the stated terminal-localization equivalence.
\end{proof}

The subcategory remains proper: the split admissible example
at the end of A-005 has a nonzero simplex-zero weight-one one-form,
whereas every object of $\mathcal C$ has the specified strict
simplex-zero identification with $S$ and hence has zero such group.
Both main questions have thus been transferred, by the actual
functors and transformations above, to this restricted class.
Their truth for that class is still undecided; no comparison
between arbitrary two of its objects or contraction of all its
zigzag loops is asserted here.
